\slajdid{07-s007}
\begin{frame}[label=07-s007]{Primer jedne PAL komponente}
\gusttekst\setlength{\parskip}{1pt}\begin{columns}[c,onlytextwidth]
\column{.61\textwidth}\centering\begin{tikzpicture}[remember picture,x=1cm,y=1cm,font=\fontsize{8}{9}\selectfont,line width=.25pt,>=Latex]
\draw ( 0.0,.10)--(0.0,-4.9);
\draw[DEplava!55] (0.0,.10)--(0.0,0.1);
\node[text=DEplava,inner sep=.2pt,fill=white] at (0.0,0.2) {0};
\draw ( 0.08,.10)--(0.08,-4.9);
\draw[DEplava!55] (0.08,.10)--(0.08,0.36);
\node[text=DEplava,inner sep=.2pt,fill=white] at (0.08,0.46) {1};
\draw ( 0.16,.10)--(0.16,-4.9);
\draw[DEplava!55] (0.16,.10)--(0.16,0.62);
\node[text=DEplava,inner sep=.2pt,fill=white] at (0.16,0.72) {2};
\draw ( 0.24,.10)--(0.24,-4.9);
\draw[DEplava!55] (0.24,.10)--(0.24,0.88);
\node[text=DEplava,inner sep=.2pt,fill=white] at (0.24,0.98) {3};
\draw ( 0.32,.10)--(0.32,-4.9);
\draw[DEplava!55] (0.32,.10)--(0.32,1.14);
\node[text=DEplava,inner sep=.2pt,fill=white] at (0.32,1.24) {4};
\draw ( 0.4,.10)--(0.4,-4.9);
\draw[DEplava!55] (0.4,.10)--(0.4,0.1);
\node[text=DEplava,inner sep=.2pt,fill=white] at (0.4,0.2) {5};
\draw ( 0.48,.10)--(0.48,-4.9);
\draw[DEplava!55] (0.48,.10)--(0.48,0.36);
\node[text=DEplava,inner sep=.2pt,fill=white] at (0.48,0.46) {6};
\draw ( 0.56,.10)--(0.56,-4.9);
\draw[DEplava!55] (0.56,.10)--(0.56,0.62);
\node[text=DEplava,inner sep=.2pt,fill=white] at (0.56,0.72) {7};
\draw ( 0.64,.10)--(0.64,-4.9);
\draw[DEplava!55] (0.64,.10)--(0.64,0.88);
\node[text=DEplava,inner sep=.2pt,fill=white] at (0.64,0.98) {8};
\draw ( 0.72,.10)--(0.72,-4.9);
\draw[DEplava!55] (0.72,.10)--(0.72,1.14);
\node[text=DEplava,inner sep=.2pt,fill=white] at (0.72,1.24) {9};
\draw ( 0.8,.10)--(0.8,-4.9);
\draw[DEplava!55] (0.8,.10)--(0.8,0.1);
\node[text=DEplava,inner sep=.2pt,fill=white] at (0.8,0.2) {10};
\draw ( 0.88,.10)--(0.88,-4.9);
\draw[DEplava!55] (0.88,.10)--(0.88,0.36);
\node[text=DEplava,inner sep=.2pt,fill=white] at (0.88,0.46) {11};
\draw ( 0.96,.10)--(0.96,-4.9);
\draw[DEplava!55] (0.96,.10)--(0.96,0.62);
\node[text=DEplava,inner sep=.2pt,fill=white] at (0.96,0.72) {12};
\draw ( 1.04,.10)--(1.04,-4.9);
\draw[DEplava!55] (1.04,.10)--(1.04,0.88);
\node[text=DEplava,inner sep=.2pt,fill=white] at (1.04,0.98) {13};
\draw ( 1.12,.10)--(1.12,-4.9);
\draw[DEplava!55] (1.12,.10)--(1.12,1.14);
\node[text=DEplava,inner sep=.2pt,fill=white] at (1.12,1.24) {14};
\draw ( 1.2,.10)--(1.2,-4.9);
\draw[DEplava!55] (1.2,.10)--(1.2,0.1);
\node[text=DEplava,inner sep=.2pt,fill=white] at (1.2,0.2) {15};
\draw ( 1.28,.10)--(1.28,-4.9);
\draw[DEplava!55] (1.28,.10)--(1.28,0.36);
\node[text=DEplava,inner sep=.2pt,fill=white] at (1.28,0.46) {16};
\draw ( 1.36,.10)--(1.36,-4.9);
\draw[DEplava!55] (1.36,.10)--(1.36,0.62);
\node[text=DEplava,inner sep=.2pt,fill=white] at (1.36,0.72) {17};
\draw ( 1.44,.10)--(1.44,-4.9);
\draw[DEplava!55] (1.44,.10)--(1.44,0.88);
\node[text=DEplava,inner sep=.2pt,fill=white] at (1.44,0.98) {18};
\draw ( 1.52,.10)--(1.52,-4.9);
\draw[DEplava!55] (1.52,.10)--(1.52,1.14);
\node[text=DEplava,inner sep=.2pt,fill=white] at (1.52,1.24) {19};
\draw ( 1.6,.10)--(1.6,-4.9);
\draw[DEplava!55] (1.6,.10)--(1.6,0.1);
\node[text=DEplava,inner sep=.2pt,fill=white] at (1.6,0.2) {20};
\draw ( 1.68,.10)--(1.68,-4.9);
\draw[DEplava!55] (1.68,.10)--(1.68,0.36);
\node[text=DEplava,inner sep=.2pt,fill=white] at (1.68,0.46) {21};
\draw ( 1.76,.10)--(1.76,-4.9);
\draw[DEplava!55] (1.76,.10)--(1.76,0.62);
\node[text=DEplava,inner sep=.2pt,fill=white] at (1.76,0.72) {22};
\draw ( 1.84,.10)--(1.84,-4.9);
\draw[DEplava!55] (1.84,.10)--(1.84,0.88);
\node[text=DEplava,inner sep=.2pt,fill=white] at (1.84,0.98) {23};
\draw ( 1.92,.10)--(1.92,-4.9);
\draw[DEplava!55] (1.92,.10)--(1.92,1.14);
\node[text=DEplava,inner sep=.2pt,fill=white] at (1.92,1.24) {24};
\draw ( 2.0,.10)--(2.0,-4.9);
\draw[DEplava!55] (2.0,.10)--(2.0,0.1);
\node[text=DEplava,inner sep=.2pt,fill=white] at (2.0,0.2) {25};
\draw ( 2.08,.10)--(2.08,-4.9);
\draw[DEplava!55] (2.08,.10)--(2.08,0.36);
\node[text=DEplava,inner sep=.2pt,fill=white] at (2.08,0.46) {26};
\draw ( 2.16,.10)--(2.16,-4.9);
\draw[DEplava!55] (2.16,.10)--(2.16,0.62);
\node[text=DEplava,inner sep=.2pt,fill=white] at (2.16,0.72) {27};
\draw ( 2.24,.10)--(2.24,-4.9);
\draw[DEplava!55] (2.24,.10)--(2.24,0.88);
\node[text=DEplava,inner sep=.2pt,fill=white] at (2.24,0.98) {28};
\draw ( 2.32,.10)--(2.32,-4.9);
\draw[DEplava!55] (2.32,.10)--(2.32,1.14);
\node[text=DEplava,inner sep=.2pt,fill=white] at (2.32,1.24) {29};
\draw ( 2.4,.10)--(2.4,-4.9);
\draw[DEplava!55] (2.4,.10)--(2.4,0.1);
\node[text=DEplava,inner sep=.2pt,fill=white] at (2.4,0.2) {30};
\draw ( 2.48,.10)--(2.48,-4.9);
\draw[DEplava!55] (2.48,.10)--(2.48,0.36);
\node[text=DEplava,inner sep=.2pt,fill=white] at (2.48,0.46) {31};
\draw[DEplava!55] (-0.23,-0.0)--(-.10,-0.0);
\node[fill=white,inner sep=.2pt,text=DEplava] at (-0.23,-0.0) {0};
\draw (-.10,-0.0)--(2.55,-0.0);
\draw[fill=white] (2.55,-0.024)--(2.63,-0.024) arc(-90:90:.024)--(2.55,0.024)--cycle;
\draw (2.654,-0.0)--(3.47,-0.0)--(3.47,-0.15);
\draw[DEplava!55] (-0.61,-0.06)--(-.10,-0.06);
\node[fill=white,inner sep=.2pt,text=DEplava] at (-0.61,-0.06) {1};
\draw (-.10,-0.06)--(2.55,-0.06);
\draw[fill=white] (2.55,-0.08399999999999999)--(2.63,-0.08399999999999999) arc(-90:90:.024)--(2.55,-0.036)--cycle;
\draw (2.654,-0.06)--(2.73,-0.06)--(2.73,-0.12)--(2.87,-0.12);
\draw[DEplava!55] (-0.99,-0.12)--(-.10,-0.12);
\node[fill=white,inner sep=.2pt,text=DEplava] at (-0.99,-0.12) {2};
\draw (-.10,-0.12)--(2.55,-0.12);
\draw[fill=white] (2.55,-0.144)--(2.63,-0.144) arc(-90:90:.024)--(2.55,-0.096)--cycle;
\draw (2.654,-0.12)--(2.73,-0.12)--(2.73,-0.16)--(2.87,-0.16);
\draw[DEplava!55] (-1.37,-0.18)--(-.10,-0.18);
\node[fill=white,inner sep=.2pt,text=DEplava] at (-1.37,-0.18) {3};
\draw (-.10,-0.18)--(2.55,-0.18);
\draw[fill=white] (2.55,-0.204)--(2.63,-0.204) arc(-90:90:.024)--(2.55,-0.156)--cycle;
\draw (2.654,-0.18)--(2.73,-0.18)--(2.73,-0.2)--(2.87,-0.2);
\draw[DEplava!55] (-1.75,-0.24)--(-.10,-0.24);
\node[fill=white,inner sep=.2pt,text=DEplava] at (-1.75,-0.24) {4};
\draw (-.10,-0.24)--(2.55,-0.24);
\draw[fill=white] (2.55,-0.264)--(2.63,-0.264) arc(-90:90:.024)--(2.55,-0.216)--cycle;
\draw (2.654,-0.24)--(2.73,-0.24)--(2.73,-0.24)--(2.87,-0.24);
\draw[DEplava!55] (-0.23,-0.3)--(-.10,-0.3);
\node[fill=white,inner sep=.2pt,text=DEplava] at (-0.23,-0.3) {5};
\draw (-.10,-0.3)--(2.55,-0.3);
\draw[fill=white] (2.55,-0.324)--(2.63,-0.324) arc(-90:90:.024)--(2.55,-0.27599999999999997)--cycle;
\draw (2.654,-0.3)--(2.73,-0.3)--(2.73,-0.28)--(2.87,-0.28);
\draw[DEplava!55] (-0.61,-0.36)--(-.10,-0.36);
\node[fill=white,inner sep=.2pt,text=DEplava] at (-0.61,-0.36) {6};
\draw (-.10,-0.36)--(2.55,-0.36);
\draw[fill=white] (2.55,-0.384)--(2.63,-0.384) arc(-90:90:.024)--(2.55,-0.33599999999999997)--cycle;
\draw (2.654,-0.36)--(2.73,-0.36)--(2.73,-0.32)--(2.87,-0.32);
\draw[DEplava!55] (-0.99,-0.42)--(-.10,-0.42);
\node[fill=white,inner sep=.2pt,text=DEplava] at (-0.99,-0.42) {7};
\draw (-.10,-0.42)--(2.55,-0.42);
\draw[fill=white] (2.55,-0.444)--(2.63,-0.444) arc(-90:90:.024)--(2.55,-0.39599999999999996)--cycle;
\draw (2.654,-0.42)--(2.73,-0.42)--(2.73,-0.36)--(2.87,-0.36);
\draw[fill=white] (2.82,-0.42) .. controls (3.02,-0.42) and (3.16,-0.39) .. (3.28,-0.24) .. controls (3.16,-0.09) and (3.02,-0.06) .. (2.82,-0.06) .. controls (2.94,-0.19999999999999998) and (2.94,-0.27999999999999997) .. cycle;
\draw (3.28,-0.24)--(3.38,-0.24);
\draw[fill=white] (3.38,-0.32999999999999996)--(3.61,-0.24)--(3.38,-0.15)--cycle;\draw[fill=white] (3.64,-0.24) circle(.03);
\draw (3.67,-0.24)--(4.10,-0.24);
\node[anchor=west,inner sep=1pt] at (4.1,-0.24) {O1};
\node[anchor=south,text=DEplava,inner sep=0pt] at (3.95,-0.22) {(19)};
\draw (-2.25,-0.51)--(-2.02,-0.51);
\draw[fill=white] (-2.02,-0.6)--(-1.80,-0.51)--(-2.02,-0.42000000000000004)--cycle;
\draw[fill=white] (-1.80,-0.5650000000000001) circle(.025);
\draw (-1.84,-0.455)--(0.0,-0.455);
\draw (-1.775,-0.5650000000000001)--(0.08,-0.5650000000000001);
\fill (0.0,-0.455) circle(.012);\fill (0.08,-0.5650000000000001) circle(.012);
\node[anchor=east,inner sep=1pt] at (-2.25,-0.51) {I2};
\node[anchor=south,text=DEplava,inner sep=0pt] at (-2.30,-0.45) {(2)};
\draw[fill=white] (3.63,-0.6)--(3.40,-0.51)--(3.63,-0.42000000000000004)--cycle;
\draw[fill=white] (3.40,-0.5650000000000001) circle(.025);
\draw (3.44,-0.455)--(0.16,-0.455);
\draw (3.375,-0.5650000000000001)--(0.24,-0.5650000000000001);
\fill (0.16,-0.455) circle(.012);\fill (0.24,-0.5650000000000001) circle(.012);
\draw (-2.25,1.55)--(4.35,1.55)--(4.35,-0.51)--(3.63,-0.51);\node[anchor=east] at (-2.25,1.55) {I1};\node[anchor=south,text=DEplava] at (-2.3,1.57) {(1)};
\draw[DEplava!55] (-1.37,-0.6)--(-.10,-0.6);
\node[fill=white,inner sep=.2pt,text=DEplava] at (-1.37,-0.6) {8};
\draw (-.10,-0.6)--(2.55,-0.6);
\draw[fill=white] (2.55,-0.624)--(2.63,-0.624) arc(-90:90:.024)--(2.55,-0.576)--cycle;
\draw (2.654,-0.6)--(3.47,-0.6)--(3.47,-0.75);
\draw[DEplava!55] (-1.75,-0.6599999999999999)--(-.10,-0.6599999999999999);
\node[fill=white,inner sep=.2pt,text=DEplava] at (-1.75,-0.6599999999999999) {9};
\draw (-.10,-0.6599999999999999)--(2.55,-0.6599999999999999);
\draw[fill=white] (2.55,-0.6839999999999999)--(2.63,-0.6839999999999999) arc(-90:90:.024)--(2.55,-0.6359999999999999)--cycle;
\draw (2.654,-0.6599999999999999)--(2.73,-0.6599999999999999)--(2.73,-0.72)--(2.87,-0.72);
\draw[DEplava!55] (-0.23,-0.72)--(-.10,-0.72);
\node[fill=white,inner sep=.2pt,text=DEplava] at (-0.23,-0.72) {10};
\draw (-.10,-0.72)--(2.55,-0.72);
\draw[fill=white] (2.55,-0.744)--(2.63,-0.744) arc(-90:90:.024)--(2.55,-0.696)--cycle;
\draw (2.654,-0.72)--(2.73,-0.72)--(2.73,-0.76)--(2.87,-0.76);
\draw[DEplava!55] (-0.61,-0.78)--(-.10,-0.78);
\node[fill=white,inner sep=.2pt,text=DEplava] at (-0.61,-0.78) {11};
\draw (-.10,-0.78)--(2.55,-0.78);
\draw[fill=white] (2.55,-0.804)--(2.63,-0.804) arc(-90:90:.024)--(2.55,-0.756)--cycle;
\draw (2.654,-0.78)--(2.73,-0.78)--(2.73,-0.7999999999999999)--(2.87,-0.7999999999999999);
\draw[DEplava!55] (-0.99,-0.84)--(-.10,-0.84);
\node[fill=white,inner sep=.2pt,text=DEplava] at (-0.99,-0.84) {12};
\draw (-.10,-0.84)--(2.55,-0.84);
\draw[fill=white] (2.55,-0.864)--(2.63,-0.864) arc(-90:90:.024)--(2.55,-0.816)--cycle;
\draw (2.654,-0.84)--(2.73,-0.84)--(2.73,-0.84)--(2.87,-0.84);
\draw[DEplava!55] (-1.37,-0.8999999999999999)--(-.10,-0.8999999999999999);
\node[fill=white,inner sep=.2pt,text=DEplava] at (-1.37,-0.8999999999999999) {13};
\draw (-.10,-0.8999999999999999)--(2.55,-0.8999999999999999);
\draw[fill=white] (2.55,-0.9239999999999999)--(2.63,-0.9239999999999999) arc(-90:90:.024)--(2.55,-0.8759999999999999)--cycle;
\draw (2.654,-0.8999999999999999)--(2.73,-0.8999999999999999)--(2.73,-0.88)--(2.87,-0.88);
\draw[DEplava!55] (-1.75,-0.96)--(-.10,-0.96);
\node[fill=white,inner sep=.2pt,text=DEplava] at (-1.75,-0.96) {14};
\draw (-.10,-0.96)--(2.55,-0.96);
\draw[fill=white] (2.55,-0.984)--(2.63,-0.984) arc(-90:90:.024)--(2.55,-0.9359999999999999)--cycle;
\draw (2.654,-0.96)--(2.73,-0.96)--(2.73,-0.9199999999999999)--(2.87,-0.9199999999999999);
\draw[DEplava!55] (-0.23,-1.02)--(-.10,-1.02);
\node[fill=white,inner sep=.2pt,text=DEplava] at (-0.23,-1.02) {15};
\draw (-.10,-1.02)--(2.55,-1.02);
\draw[fill=white] (2.55,-1.044)--(2.63,-1.044) arc(-90:90:.024)--(2.55,-0.996)--cycle;
\draw (2.654,-1.02)--(2.73,-1.02)--(2.73,-0.96)--(2.87,-0.96);
\draw[fill=white] (2.82,-1.02) .. controls (3.02,-1.02) and (3.16,-0.99) .. (3.28,-0.84) .. controls (3.16,-0.69) and (3.02,-0.6599999999999999) .. (2.82,-0.6599999999999999) .. controls (2.94,-0.7999999999999999) and (2.94,-0.88) .. cycle;
\draw (3.28,-0.84)--(3.38,-0.84);
\draw[fill=white] (3.38,-0.9299999999999999)--(3.61,-0.84)--(3.38,-0.75)--cycle;\draw[fill=white] (3.64,-0.84) circle(.03);
\draw (3.67,-0.84)--(4.10,-0.84);
\node[anchor=west,inner sep=1pt] at (4.1,-0.84) {IO2};
\node[anchor=south,text=DEplava,inner sep=0pt] at (3.95,-0.82) {(18)};
\draw (-2.25,-1.1099999999999999)--(-2.02,-1.1099999999999999);
\draw[fill=white] (-2.02,-1.2)--(-1.80,-1.1099999999999999)--(-2.02,-1.0199999999999998)--cycle;
\draw[fill=white] (-1.80,-1.1649999999999998) circle(.025);
\draw (-1.84,-1.055)--(0.32,-1.055);
\draw (-1.775,-1.1649999999999998)--(0.4,-1.1649999999999998);
\fill (0.32,-1.055) circle(.012);\fill (0.4,-1.1649999999999998) circle(.012);
\node[anchor=east,inner sep=1pt] at (-2.25,-1.1099999999999999) {I3};
\node[anchor=south,text=DEplava,inner sep=0pt] at (-2.30,-1.0499999999999998) {(3)};
\draw[fill=white] (3.63,-1.2)--(3.40,-1.1099999999999999)--(3.63,-1.0199999999999998)--cycle;
\draw[fill=white] (3.40,-1.1649999999999998) circle(.025);
\draw (3.44,-1.055)--(0.48,-1.055);
\draw (3.375,-1.1649999999999998)--(0.56,-1.1649999999999998);
\fill (0.48,-1.055) circle(.012);\fill (0.56,-1.1649999999999998) circle(.012);
\draw (3.85,-0.84)--(3.85,-1.1099999999999999)--(3.63,-1.1099999999999999);\fill (3.85,-0.84) circle(.012);
\draw[DEplava!55] (-0.61,-1.2)--(-.10,-1.2);
\node[fill=white,inner sep=.2pt,text=DEplava] at (-0.61,-1.2) {16};
\draw (-.10,-1.2)--(2.55,-1.2);
\draw[fill=white] (2.55,-1.224)--(2.63,-1.224) arc(-90:90:.024)--(2.55,-1.176)--cycle;
\draw (2.654,-1.2)--(3.47,-1.2)--(3.47,-1.3499999999999999);
\draw[DEplava!55] (-0.99,-1.26)--(-.10,-1.26);
\node[fill=white,inner sep=.2pt,text=DEplava] at (-0.99,-1.26) {17};
\draw (-.10,-1.26)--(2.55,-1.26);
\draw[fill=white] (2.55,-1.284)--(2.63,-1.284) arc(-90:90:.024)--(2.55,-1.236)--cycle;
\draw (2.654,-1.26)--(2.73,-1.26)--(2.73,-1.3199999999999998)--(2.87,-1.3199999999999998);
\draw[DEplava!55] (-1.37,-1.3199999999999998)--(-.10,-1.3199999999999998);
\node[fill=white,inner sep=.2pt,text=DEplava] at (-1.37,-1.3199999999999998) {18};
\draw (-.10,-1.3199999999999998)--(2.55,-1.3199999999999998);
\draw[fill=white] (2.55,-1.3439999999999999)--(2.63,-1.3439999999999999) arc(-90:90:.024)--(2.55,-1.2959999999999998)--cycle;
\draw (2.654,-1.3199999999999998)--(2.73,-1.3199999999999998)--(2.73,-1.3599999999999999)--(2.87,-1.3599999999999999);
\draw[DEplava!55] (-1.75,-1.38)--(-.10,-1.38);
\node[fill=white,inner sep=.2pt,text=DEplava] at (-1.75,-1.38) {19};
\draw (-.10,-1.38)--(2.55,-1.38);
\draw[fill=white] (2.55,-1.404)--(2.63,-1.404) arc(-90:90:.024)--(2.55,-1.3559999999999999)--cycle;
\draw (2.654,-1.38)--(2.73,-1.38)--(2.73,-1.4)--(2.87,-1.4);
\draw[DEplava!55] (-0.23,-1.44)--(-.10,-1.44);
\node[fill=white,inner sep=.2pt,text=DEplava] at (-0.23,-1.44) {20};
\draw (-.10,-1.44)--(2.55,-1.44);
\draw[fill=white] (2.55,-1.464)--(2.63,-1.464) arc(-90:90:.024)--(2.55,-1.416)--cycle;
\draw (2.654,-1.44)--(2.73,-1.44)--(2.73,-1.44)--(2.87,-1.44);
\draw[DEplava!55] (-0.61,-1.5)--(-.10,-1.5);
\node[fill=white,inner sep=.2pt,text=DEplava] at (-0.61,-1.5) {21};
\draw (-.10,-1.5)--(2.55,-1.5);
\draw[fill=white] (2.55,-1.524)--(2.63,-1.524) arc(-90:90:.024)--(2.55,-1.476)--cycle;
\draw (2.654,-1.5)--(2.73,-1.5)--(2.73,-1.4799999999999998)--(2.87,-1.4799999999999998);
\draw[DEplava!55] (-0.99,-1.56)--(-.10,-1.56);
\node[fill=white,inner sep=.2pt,text=DEplava] at (-0.99,-1.56) {22};
\draw (-.10,-1.56)--(2.55,-1.56);
\draw[fill=white] (2.55,-1.584)--(2.63,-1.584) arc(-90:90:.024)--(2.55,-1.536)--cycle;
\draw (2.654,-1.56)--(2.73,-1.56)--(2.73,-1.5199999999999998)--(2.87,-1.5199999999999998);
\draw[DEplava!55] (-1.37,-1.6199999999999999)--(-.10,-1.6199999999999999);
\node[fill=white,inner sep=.2pt,text=DEplava] at (-1.37,-1.6199999999999999) {23};
\draw (-.10,-1.6199999999999999)--(2.55,-1.6199999999999999);
\draw[fill=white] (2.55,-1.644)--(2.63,-1.644) arc(-90:90:.024)--(2.55,-1.5959999999999999)--cycle;
\draw (2.654,-1.6199999999999999)--(2.73,-1.6199999999999999)--(2.73,-1.5599999999999998)--(2.87,-1.5599999999999998);
\draw[fill=white] (2.82,-1.6199999999999999) .. controls (3.02,-1.6199999999999999) and (3.16,-1.5899999999999999) .. (3.28,-1.44) .. controls (3.16,-1.29) and (3.02,-1.26) .. (2.82,-1.26) .. controls (2.94,-1.4) and (2.94,-1.48) .. cycle;
\draw (3.28,-1.44)--(3.38,-1.44);
\draw[fill=white] (3.38,-1.53)--(3.61,-1.44)--(3.38,-1.3499999999999999)--cycle;\draw[fill=white] (3.64,-1.44) circle(.03);
\draw (3.67,-1.44)--(4.10,-1.44);
\node[anchor=west,inner sep=1pt] at (4.1,-1.44) {IO3};
\node[anchor=south,text=DEplava,inner sep=0pt] at (3.95,-1.42) {(17)};
\draw (-2.25,-1.71)--(-2.02,-1.71);
\draw[fill=white] (-2.02,-1.8)--(-1.80,-1.71)--(-2.02,-1.6199999999999999)--cycle;
\draw[fill=white] (-1.80,-1.765) circle(.025);
\draw (-1.84,-1.655)--(0.64,-1.655);
\draw (-1.775,-1.765)--(0.72,-1.765);
\fill (0.64,-1.655) circle(.012);\fill (0.72,-1.765) circle(.012);
\node[anchor=east,inner sep=1pt] at (-2.25,-1.71) {I4};
\node[anchor=south,text=DEplava,inner sep=0pt] at (-2.30,-1.65) {(4)};
\draw[fill=white] (3.63,-1.8)--(3.40,-1.71)--(3.63,-1.6199999999999999)--cycle;
\draw[fill=white] (3.40,-1.765) circle(.025);
\draw (3.44,-1.655)--(0.8,-1.655);
\draw (3.375,-1.765)--(0.88,-1.765);
\fill (0.8,-1.655) circle(.012);\fill (0.88,-1.765) circle(.012);
\draw (3.85,-1.44)--(3.85,-1.71)--(3.63,-1.71);\fill (3.85,-1.44) circle(.012);
\draw[DEplava!55] (-1.75,-1.7999999999999998)--(-.10,-1.7999999999999998);
\node[fill=white,inner sep=.2pt,text=DEplava] at (-1.75,-1.7999999999999998) {24};
\draw (-.10,-1.7999999999999998)--(2.55,-1.7999999999999998);
\draw[fill=white] (2.55,-1.8239999999999998)--(2.63,-1.8239999999999998) arc(-90:90:.024)--(2.55,-1.7759999999999998)--cycle;
\draw (2.654,-1.7999999999999998)--(3.47,-1.7999999999999998)--(3.47,-1.95);
\draw[DEplava!55] (-0.23,-1.8599999999999999)--(-.10,-1.8599999999999999);
\node[fill=white,inner sep=.2pt,text=DEplava] at (-0.23,-1.8599999999999999) {25};
\draw (-.10,-1.8599999999999999)--(2.55,-1.8599999999999999);
\draw[fill=white] (2.55,-1.884)--(2.63,-1.884) arc(-90:90:.024)--(2.55,-1.8359999999999999)--cycle;
\draw (2.654,-1.8599999999999999)--(2.73,-1.8599999999999999)--(2.73,-1.92)--(2.87,-1.92);
\draw[DEplava!55] (-0.61,-1.92)--(-.10,-1.92);
\node[fill=white,inner sep=.2pt,text=DEplava] at (-0.61,-1.92) {26};
\draw (-.10,-1.92)--(2.55,-1.92);
\draw[fill=white] (2.55,-1.944)--(2.63,-1.944) arc(-90:90:.024)--(2.55,-1.896)--cycle;
\draw (2.654,-1.92)--(2.73,-1.92)--(2.73,-1.96)--(2.87,-1.96);
\draw[DEplava!55] (-0.99,-1.9799999999999998)--(-.10,-1.9799999999999998);
\node[fill=white,inner sep=.2pt,text=DEplava] at (-0.99,-1.9799999999999998) {27};
\draw (-.10,-1.9799999999999998)--(2.55,-1.9799999999999998);
\draw[fill=white] (2.55,-2.0039999999999996)--(2.63,-2.0039999999999996) arc(-90:90:.024)--(2.55,-1.9559999999999997)--cycle;
\draw (2.654,-1.9799999999999998)--(2.73,-1.9799999999999998)--(2.73,-2.0)--(2.87,-2.0);
\draw[DEplava!55] (-1.37,-2.04)--(-.10,-2.04);
\node[fill=white,inner sep=.2pt,text=DEplava] at (-1.37,-2.04) {28};
\draw (-.10,-2.04)--(2.55,-2.04);
\draw[fill=white] (2.55,-2.064)--(2.63,-2.064) arc(-90:90:.024)--(2.55,-2.016)--cycle;
\draw (2.654,-2.04)--(2.73,-2.04)--(2.73,-2.04)--(2.87,-2.04);
\draw[DEplava!55] (-1.75,-2.0999999999999996)--(-.10,-2.0999999999999996);
\node[fill=white,inner sep=.2pt,text=DEplava] at (-1.75,-2.0999999999999996) {29};
\draw (-.10,-2.0999999999999996)--(2.55,-2.0999999999999996);
\draw[fill=white] (2.55,-2.1239999999999997)--(2.63,-2.1239999999999997) arc(-90:90:.024)--(2.55,-2.0759999999999996)--cycle;
\draw (2.654,-2.0999999999999996)--(2.73,-2.0999999999999996)--(2.73,-2.08)--(2.87,-2.08);
\draw[DEplava!55] (-0.23,-2.1599999999999997)--(-.10,-2.1599999999999997);
\node[fill=white,inner sep=.2pt,text=DEplava] at (-0.23,-2.1599999999999997) {30};
\draw (-.10,-2.1599999999999997)--(2.55,-2.1599999999999997);
\draw[fill=white] (2.55,-2.1839999999999997)--(2.63,-2.1839999999999997) arc(-90:90:.024)--(2.55,-2.1359999999999997)--cycle;
\draw (2.654,-2.1599999999999997)--(2.73,-2.1599999999999997)--(2.73,-2.12)--(2.87,-2.12);
\draw[DEplava!55] (-0.61,-2.2199999999999998)--(-.10,-2.2199999999999998);
\node[fill=white,inner sep=.2pt,text=DEplava] at (-0.61,-2.2199999999999998) {31};
\draw (-.10,-2.2199999999999998)--(2.55,-2.2199999999999998);
\draw[fill=white] (2.55,-2.2439999999999998)--(2.63,-2.2439999999999998) arc(-90:90:.024)--(2.55,-2.1959999999999997)--cycle;
\draw (2.654,-2.2199999999999998)--(2.73,-2.2199999999999998)--(2.73,-2.16)--(2.87,-2.16);
\draw[fill=white] (2.82,-2.22) .. controls (3.02,-2.22) and (3.16,-2.19) .. (3.28,-2.04) .. controls (3.16,-1.8900000000000001) and (3.02,-1.86) .. (2.82,-1.86) .. controls (2.94,-2.0) and (2.94,-2.08) .. cycle;
\draw (3.28,-2.04)--(3.38,-2.04);
\draw[fill=white] (3.38,-2.13)--(3.61,-2.04)--(3.38,-1.95)--cycle;\draw[fill=white] (3.64,-2.04) circle(.03);
\draw (3.67,-2.04)--(4.10,-2.04);
\node[anchor=west,inner sep=1pt] at (4.1,-2.04) {IO4};
\node[anchor=south,text=DEplava,inner sep=0pt] at (3.95,-2.02) {(16)};
\draw (-2.25,-2.3099999999999996)--(-2.02,-2.3099999999999996);
\draw[fill=white] (-2.02,-2.3999999999999995)--(-1.80,-2.3099999999999996)--(-2.02,-2.2199999999999998)--cycle;
\draw[fill=white] (-1.80,-2.3649999999999998) circle(.025);
\draw (-1.84,-2.2549999999999994)--(0.96,-2.2549999999999994);
\draw (-1.775,-2.3649999999999998)--(1.04,-2.3649999999999998);
\fill (0.96,-2.2549999999999994) circle(.012);\fill (1.04,-2.3649999999999998) circle(.012);
\node[anchor=east,inner sep=1pt] at (-2.25,-2.3099999999999996) {I5};
\node[anchor=south,text=DEplava,inner sep=0pt] at (-2.30,-2.2499999999999996) {(5)};
\draw[fill=white] (3.63,-2.3999999999999995)--(3.40,-2.3099999999999996)--(3.63,-2.2199999999999998)--cycle;
\draw[fill=white] (3.40,-2.3649999999999998) circle(.025);
\draw (3.44,-2.2549999999999994)--(1.12,-2.2549999999999994);
\draw (3.375,-2.3649999999999998)--(1.2,-2.3649999999999998);
\fill (1.12,-2.2549999999999994) circle(.012);\fill (1.2,-2.3649999999999998) circle(.012);
\draw (3.85,-2.04)--(3.85,-2.3099999999999996)--(3.63,-2.3099999999999996);\fill (3.85,-2.04) circle(.012);
\draw[DEplava!55] (-0.99,-2.4)--(-.10,-2.4);
\node[fill=white,inner sep=.2pt,text=DEplava] at (-0.99,-2.4) {32};
\draw (-.10,-2.4)--(2.55,-2.4);
\draw[fill=white] (2.55,-2.424)--(2.63,-2.424) arc(-90:90:.024)--(2.55,-2.376)--cycle;
\draw (2.654,-2.4)--(3.47,-2.4)--(3.47,-2.55);
\draw[DEplava!55] (-1.37,-2.46)--(-.10,-2.46);
\node[fill=white,inner sep=.2pt,text=DEplava] at (-1.37,-2.46) {33};
\draw (-.10,-2.46)--(2.55,-2.46);
\draw[fill=white] (2.55,-2.484)--(2.63,-2.484) arc(-90:90:.024)--(2.55,-2.436)--cycle;
\draw (2.654,-2.46)--(2.73,-2.46)--(2.73,-2.5199999999999996)--(2.87,-2.5199999999999996);
\draw[DEplava!55] (-1.75,-2.52)--(-.10,-2.52);
\node[fill=white,inner sep=.2pt,text=DEplava] at (-1.75,-2.52) {34};
\draw (-.10,-2.52)--(2.55,-2.52);
\draw[fill=white] (2.55,-2.544)--(2.63,-2.544) arc(-90:90:.024)--(2.55,-2.496)--cycle;
\draw (2.654,-2.52)--(2.73,-2.52)--(2.73,-2.5599999999999996)--(2.87,-2.5599999999999996);
\draw[DEplava!55] (-0.23,-2.58)--(-.10,-2.58);
\node[fill=white,inner sep=.2pt,text=DEplava] at (-0.23,-2.58) {35};
\draw (-.10,-2.58)--(2.55,-2.58);
\draw[fill=white] (2.55,-2.604)--(2.63,-2.604) arc(-90:90:.024)--(2.55,-2.556)--cycle;
\draw (2.654,-2.58)--(2.73,-2.58)--(2.73,-2.5999999999999996)--(2.87,-2.5999999999999996);
\draw[DEplava!55] (-0.61,-2.6399999999999997)--(-.10,-2.6399999999999997);
\node[fill=white,inner sep=.2pt,text=DEplava] at (-0.61,-2.6399999999999997) {36};
\draw (-.10,-2.6399999999999997)--(2.55,-2.6399999999999997);
\draw[fill=white] (2.55,-2.6639999999999997)--(2.63,-2.6639999999999997) arc(-90:90:.024)--(2.55,-2.6159999999999997)--cycle;
\draw (2.654,-2.6399999999999997)--(2.73,-2.6399999999999997)--(2.73,-2.6399999999999997)--(2.87,-2.6399999999999997);
\draw[DEplava!55] (-0.99,-2.6999999999999997)--(-.10,-2.6999999999999997);
\node[fill=white,inner sep=.2pt,text=DEplava] at (-0.99,-2.6999999999999997) {37};
\draw (-.10,-2.6999999999999997)--(2.55,-2.6999999999999997);
\draw[fill=white] (2.55,-2.7239999999999998)--(2.63,-2.7239999999999998) arc(-90:90:.024)--(2.55,-2.6759999999999997)--cycle;
\draw (2.654,-2.6999999999999997)--(2.73,-2.6999999999999997)--(2.73,-2.6799999999999997)--(2.87,-2.6799999999999997);
\draw[DEplava!55] (-1.37,-2.76)--(-.10,-2.76);
\node[fill=white,inner sep=.2pt,text=DEplava] at (-1.37,-2.76) {38};
\draw (-.10,-2.76)--(2.55,-2.76);
\draw[fill=white] (2.55,-2.784)--(2.63,-2.784) arc(-90:90:.024)--(2.55,-2.7359999999999998)--cycle;
\draw (2.654,-2.76)--(2.73,-2.76)--(2.73,-2.7199999999999998)--(2.87,-2.7199999999999998);
\draw[DEplava!55] (-1.75,-2.82)--(-.10,-2.82);
\node[fill=white,inner sep=.2pt,text=DEplava] at (-1.75,-2.82) {39};
\draw (-.10,-2.82)--(2.55,-2.82);
\draw[fill=white] (2.55,-2.844)--(2.63,-2.844) arc(-90:90:.024)--(2.55,-2.796)--cycle;
\draw (2.654,-2.82)--(2.73,-2.82)--(2.73,-2.76)--(2.87,-2.76);
\draw[fill=white] (2.82,-2.82) .. controls (3.02,-2.82) and (3.16,-2.7899999999999996) .. (3.28,-2.6399999999999997) .. controls (3.16,-2.4899999999999998) and (3.02,-2.4599999999999995) .. (2.82,-2.4599999999999995) .. controls (2.94,-2.5999999999999996) and (2.94,-2.6799999999999997) .. cycle;
\draw (3.28,-2.6399999999999997)--(3.38,-2.6399999999999997);
\draw[fill=white] (3.38,-2.7299999999999995)--(3.61,-2.6399999999999997)--(3.38,-2.55)--cycle;\draw[fill=white] (3.64,-2.6399999999999997) circle(.03);
\draw (3.67,-2.6399999999999997)--(4.10,-2.6399999999999997);
\node[anchor=west,inner sep=1pt] at (4.1,-2.6399999999999997) {IO5};
\node[anchor=south,text=DEplava,inner sep=0pt] at (3.95,-2.6199999999999997) {(15)};
\draw (-2.25,-2.91)--(-2.02,-2.91);
\draw[fill=white] (-2.02,-3.0)--(-1.80,-2.91)--(-2.02,-2.8200000000000003)--cycle;
\draw[fill=white] (-1.80,-2.9650000000000003) circle(.025);
\draw (-1.84,-2.855)--(1.28,-2.855);
\draw (-1.775,-2.9650000000000003)--(1.36,-2.9650000000000003);
\fill (1.28,-2.855) circle(.012);\fill (1.36,-2.9650000000000003) circle(.012);
\node[anchor=east,inner sep=1pt] at (-2.25,-2.91) {I6};
\node[anchor=south,text=DEplava,inner sep=0pt] at (-2.30,-2.85) {(6)};
\draw[fill=white] (3.63,-3.0)--(3.40,-2.91)--(3.63,-2.8200000000000003)--cycle;
\draw[fill=white] (3.40,-2.9650000000000003) circle(.025);
\draw (3.44,-2.855)--(1.44,-2.855);
\draw (3.375,-2.9650000000000003)--(1.52,-2.9650000000000003);
\fill (1.44,-2.855) circle(.012);\fill (1.52,-2.9650000000000003) circle(.012);
\draw (3.85,-2.6399999999999997)--(3.85,-2.91)--(3.63,-2.91);\fill (3.85,-2.6399999999999997) circle(.012);
\draw[DEplava!55] (-0.23,-3.0)--(-.10,-3.0);
\node[fill=white,inner sep=.2pt,text=DEplava] at (-0.23,-3.0) {40};
\draw (-.10,-3.0)--(2.55,-3.0);
\draw[fill=white] (2.55,-3.024)--(2.63,-3.024) arc(-90:90:.024)--(2.55,-2.976)--cycle;
\draw (2.654,-3.0)--(3.47,-3.0)--(3.47,-3.1500000000000004);
\draw[DEplava!55] (-0.61,-3.06)--(-.10,-3.06);
\node[fill=white,inner sep=.2pt,text=DEplava] at (-0.61,-3.06) {41};
\draw (-.10,-3.06)--(2.55,-3.06);
\draw[fill=white] (2.55,-3.084)--(2.63,-3.084) arc(-90:90:.024)--(2.55,-3.036)--cycle;
\draw (2.654,-3.06)--(2.73,-3.06)--(2.73,-3.12)--(2.87,-3.12);
\draw[DEplava!55] (-0.99,-3.12)--(-.10,-3.12);
\node[fill=white,inner sep=.2pt,text=DEplava] at (-0.99,-3.12) {42};
\draw (-.10,-3.12)--(2.55,-3.12);
\draw[fill=white] (2.55,-3.144)--(2.63,-3.144) arc(-90:90:.024)--(2.55,-3.096)--cycle;
\draw (2.654,-3.12)--(2.73,-3.12)--(2.73,-3.16)--(2.87,-3.16);
\draw[DEplava!55] (-1.37,-3.18)--(-.10,-3.18);
\node[fill=white,inner sep=.2pt,text=DEplava] at (-1.37,-3.18) {43};
\draw (-.10,-3.18)--(2.55,-3.18);
\draw[fill=white] (2.55,-3.204)--(2.63,-3.204) arc(-90:90:.024)--(2.55,-3.156)--cycle;
\draw (2.654,-3.18)--(2.73,-3.18)--(2.73,-3.2)--(2.87,-3.2);
\draw[DEplava!55] (-1.75,-3.24)--(-.10,-3.24);
\node[fill=white,inner sep=.2pt,text=DEplava] at (-1.75,-3.24) {44};
\draw (-.10,-3.24)--(2.55,-3.24);
\draw[fill=white] (2.55,-3.2640000000000002)--(2.63,-3.2640000000000002) arc(-90:90:.024)--(2.55,-3.216)--cycle;
\draw (2.654,-3.24)--(2.73,-3.24)--(2.73,-3.24)--(2.87,-3.24);
\draw[DEplava!55] (-0.23,-3.3)--(-.10,-3.3);
\node[fill=white,inner sep=.2pt,text=DEplava] at (-0.23,-3.3) {45};
\draw (-.10,-3.3)--(2.55,-3.3);
\draw[fill=white] (2.55,-3.324)--(2.63,-3.324) arc(-90:90:.024)--(2.55,-3.276)--cycle;
\draw (2.654,-3.3)--(2.73,-3.3)--(2.73,-3.2800000000000002)--(2.87,-3.2800000000000002);
\draw[DEplava!55] (-0.61,-3.36)--(-.10,-3.36);
\node[fill=white,inner sep=.2pt,text=DEplava] at (-0.61,-3.36) {46};
\draw (-.10,-3.36)--(2.55,-3.36);
\draw[fill=white] (2.55,-3.384)--(2.63,-3.384) arc(-90:90:.024)--(2.55,-3.336)--cycle;
\draw (2.654,-3.36)--(2.73,-3.36)--(2.73,-3.3200000000000003)--(2.87,-3.3200000000000003);
\draw[DEplava!55] (-0.99,-3.42)--(-.10,-3.42);
\node[fill=white,inner sep=.2pt,text=DEplava] at (-0.99,-3.42) {47};
\draw (-.10,-3.42)--(2.55,-3.42);
\draw[fill=white] (2.55,-3.444)--(2.63,-3.444) arc(-90:90:.024)--(2.55,-3.396)--cycle;
\draw (2.654,-3.42)--(2.73,-3.42)--(2.73,-3.3600000000000003)--(2.87,-3.3600000000000003);
\draw[fill=white] (2.82,-3.4200000000000004) .. controls (3.02,-3.4200000000000004) and (3.16,-3.39) .. (3.28,-3.24) .. controls (3.16,-3.0900000000000003) and (3.02,-3.06) .. (2.82,-3.06) .. controls (2.94,-3.2) and (2.94,-3.2800000000000002) .. cycle;
\draw (3.28,-3.24)--(3.38,-3.24);
\draw[fill=white] (3.38,-3.33)--(3.61,-3.24)--(3.38,-3.1500000000000004)--cycle;\draw[fill=white] (3.64,-3.24) circle(.03);
\draw (3.67,-3.24)--(4.10,-3.24);
\node[anchor=west,inner sep=1pt] at (4.1,-3.24) {IO6};
\node[anchor=south,text=DEplava,inner sep=0pt] at (3.95,-3.22) {(14)};
\draw (-2.25,-3.51)--(-2.02,-3.51);
\draw[fill=white] (-2.02,-3.5999999999999996)--(-1.80,-3.51)--(-2.02,-3.42)--cycle;
\draw[fill=white] (-1.80,-3.565) circle(.025);
\draw (-1.84,-3.4549999999999996)--(1.6,-3.4549999999999996);
\draw (-1.775,-3.565)--(1.68,-3.565);
\fill (1.6,-3.4549999999999996) circle(.012);\fill (1.68,-3.565) circle(.012);
\node[anchor=east,inner sep=1pt] at (-2.25,-3.51) {I7};
\node[anchor=south,text=DEplava,inner sep=0pt] at (-2.30,-3.4499999999999997) {(7)};
\draw[fill=white] (3.63,-3.5999999999999996)--(3.40,-3.51)--(3.63,-3.42)--cycle;
\draw[fill=white] (3.40,-3.565) circle(.025);
\draw (3.44,-3.4549999999999996)--(1.76,-3.4549999999999996);
\draw (3.375,-3.565)--(1.84,-3.565);
\fill (1.76,-3.4549999999999996) circle(.012);\fill (1.84,-3.565) circle(.012);
\draw (3.85,-3.24)--(3.85,-3.51)--(3.63,-3.51);\fill (3.85,-3.24) circle(.012);
\draw[DEplava!55] (-1.37,-3.5999999999999996)--(-.10,-3.5999999999999996);
\node[fill=white,inner sep=.2pt,text=DEplava] at (-1.37,-3.5999999999999996) {48};
\draw (-.10,-3.5999999999999996)--(2.55,-3.5999999999999996);
\draw[fill=white] (2.55,-3.6239999999999997)--(2.63,-3.6239999999999997) arc(-90:90:.024)--(2.55,-3.5759999999999996)--cycle;
\draw (2.654,-3.5999999999999996)--(3.47,-3.5999999999999996)--(3.47,-3.75);
\draw[DEplava!55] (-1.75,-3.6599999999999997)--(-.10,-3.6599999999999997);
\node[fill=white,inner sep=.2pt,text=DEplava] at (-1.75,-3.6599999999999997) {49};
\draw (-.10,-3.6599999999999997)--(2.55,-3.6599999999999997);
\draw[fill=white] (2.55,-3.6839999999999997)--(2.63,-3.6839999999999997) arc(-90:90:.024)--(2.55,-3.6359999999999997)--cycle;
\draw (2.654,-3.6599999999999997)--(2.73,-3.6599999999999997)--(2.73,-3.7199999999999998)--(2.87,-3.7199999999999998);
\draw[DEplava!55] (-0.23,-3.7199999999999998)--(-.10,-3.7199999999999998);
\node[fill=white,inner sep=.2pt,text=DEplava] at (-0.23,-3.7199999999999998) {50};
\draw (-.10,-3.7199999999999998)--(2.55,-3.7199999999999998);
\draw[fill=white] (2.55,-3.7439999999999998)--(2.63,-3.7439999999999998) arc(-90:90:.024)--(2.55,-3.6959999999999997)--cycle;
\draw (2.654,-3.7199999999999998)--(2.73,-3.7199999999999998)--(2.73,-3.76)--(2.87,-3.76);
\draw[DEplava!55] (-0.61,-3.78)--(-.10,-3.78);
\node[fill=white,inner sep=.2pt,text=DEplava] at (-0.61,-3.78) {51};
\draw (-.10,-3.78)--(2.55,-3.78);
\draw[fill=white] (2.55,-3.804)--(2.63,-3.804) arc(-90:90:.024)--(2.55,-3.756)--cycle;
\draw (2.654,-3.78)--(2.73,-3.78)--(2.73,-3.8)--(2.87,-3.8);
\draw[DEplava!55] (-0.99,-3.84)--(-.10,-3.84);
\node[fill=white,inner sep=.2pt,text=DEplava] at (-0.99,-3.84) {52};
\draw (-.10,-3.84)--(2.55,-3.84);
\draw[fill=white] (2.55,-3.864)--(2.63,-3.864) arc(-90:90:.024)--(2.55,-3.816)--cycle;
\draw (2.654,-3.84)--(2.73,-3.84)--(2.73,-3.84)--(2.87,-3.84);
\draw[DEplava!55] (-1.37,-3.8999999999999995)--(-.10,-3.8999999999999995);
\node[fill=white,inner sep=.2pt,text=DEplava] at (-1.37,-3.8999999999999995) {53};
\draw (-.10,-3.8999999999999995)--(2.55,-3.8999999999999995);
\draw[fill=white] (2.55,-3.9239999999999995)--(2.63,-3.9239999999999995) arc(-90:90:.024)--(2.55,-3.8759999999999994)--cycle;
\draw (2.654,-3.8999999999999995)--(2.73,-3.8999999999999995)--(2.73,-3.88)--(2.87,-3.88);
\draw[DEplava!55] (-1.75,-3.9599999999999995)--(-.10,-3.9599999999999995);
\node[fill=white,inner sep=.2pt,text=DEplava] at (-1.75,-3.9599999999999995) {54};
\draw (-.10,-3.9599999999999995)--(2.55,-3.9599999999999995);
\draw[fill=white] (2.55,-3.9839999999999995)--(2.63,-3.9839999999999995) arc(-90:90:.024)--(2.55,-3.9359999999999995)--cycle;
\draw (2.654,-3.9599999999999995)--(2.73,-3.9599999999999995)--(2.73,-3.92)--(2.87,-3.92);
\draw[DEplava!55] (-0.23,-4.02)--(-.10,-4.02);
\node[fill=white,inner sep=.2pt,text=DEplava] at (-0.23,-4.02) {55};
\draw (-.10,-4.02)--(2.55,-4.02);
\draw[fill=white] (2.55,-4.044)--(2.63,-4.044) arc(-90:90:.024)--(2.55,-3.9959999999999996)--cycle;
\draw (2.654,-4.02)--(2.73,-4.02)--(2.73,-3.96)--(2.87,-3.96);
\draw[fill=white] (2.82,-4.02) .. controls (3.02,-4.02) and (3.16,-3.9899999999999998) .. (3.28,-3.84) .. controls (3.16,-3.69) and (3.02,-3.6599999999999997) .. (2.82,-3.6599999999999997) .. controls (2.94,-3.8) and (2.94,-3.88) .. cycle;
\draw (3.28,-3.84)--(3.38,-3.84);
\draw[fill=white] (3.38,-3.9299999999999997)--(3.61,-3.84)--(3.38,-3.75)--cycle;\draw[fill=white] (3.64,-3.84) circle(.03);
\draw (3.67,-3.84)--(4.10,-3.84);
\node[anchor=west,inner sep=1pt] at (4.1,-3.84) {IO7};
\node[anchor=south,text=DEplava,inner sep=0pt] at (3.95,-3.82) {(13)};
\draw (-2.25,-4.109999999999999)--(-2.02,-4.109999999999999);
\draw[fill=white] (-2.02,-4.199999999999999)--(-1.80,-4.109999999999999)--(-2.02,-4.02)--cycle;
\draw[fill=white] (-1.80,-4.164999999999999) circle(.025);
\draw (-1.84,-4.055)--(1.92,-4.055);
\draw (-1.775,-4.164999999999999)--(2.0,-4.164999999999999);
\fill (1.92,-4.055) circle(.012);\fill (2.0,-4.164999999999999) circle(.012);
\node[anchor=east,inner sep=1pt] at (-2.25,-4.109999999999999) {I8};
\node[anchor=south,text=DEplava,inner sep=0pt] at (-2.30,-4.05) {(8)};
\draw[fill=white] (3.63,-4.199999999999999)--(3.40,-4.109999999999999)--(3.63,-4.02)--cycle;
\draw[fill=white] (3.40,-4.164999999999999) circle(.025);
\draw (3.44,-4.055)--(2.08,-4.055);
\draw (3.375,-4.164999999999999)--(2.16,-4.164999999999999);
\fill (2.08,-4.055) circle(.012);\fill (2.16,-4.164999999999999) circle(.012);
\draw (3.85,-3.84)--(3.85,-4.109999999999999)--(3.63,-4.109999999999999);\fill (3.85,-3.84) circle(.012);
\draw[DEplava!55] (-0.61,-4.2)--(-.10,-4.2);
\node[fill=white,inner sep=.2pt,text=DEplava] at (-0.61,-4.2) {56};
\draw (-.10,-4.2)--(2.55,-4.2);
\draw[fill=white] (2.55,-4.224)--(2.63,-4.224) arc(-90:90:.024)--(2.55,-4.176)--cycle;
\draw (2.654,-4.2)--(3.47,-4.2)--(3.47,-4.3500000000000005);
\draw[DEplava!55] (-0.99,-4.26)--(-.10,-4.26);
\node[fill=white,inner sep=.2pt,text=DEplava] at (-0.99,-4.26) {57};
\draw (-.10,-4.26)--(2.55,-4.26);
\draw[fill=white] (2.55,-4.284)--(2.63,-4.284) arc(-90:90:.024)--(2.55,-4.236)--cycle;
\draw (2.654,-4.26)--(2.73,-4.26)--(2.73,-4.32)--(2.87,-4.32);
\draw[DEplava!55] (-1.37,-4.32)--(-.10,-4.32);
\node[fill=white,inner sep=.2pt,text=DEplava] at (-1.37,-4.32) {58};
\draw (-.10,-4.32)--(2.55,-4.32);
\draw[fill=white] (2.55,-4.344)--(2.63,-4.344) arc(-90:90:.024)--(2.55,-4.296)--cycle;
\draw (2.654,-4.32)--(2.73,-4.32)--(2.73,-4.36)--(2.87,-4.36);
\draw[DEplava!55] (-1.75,-4.38)--(-.10,-4.38);
\node[fill=white,inner sep=.2pt,text=DEplava] at (-1.75,-4.38) {59};
\draw (-.10,-4.38)--(2.55,-4.38);
\draw[fill=white] (2.55,-4.404)--(2.63,-4.404) arc(-90:90:.024)--(2.55,-4.356)--cycle;
\draw (2.654,-4.38)--(2.73,-4.38)--(2.73,-4.4)--(2.87,-4.4);
\draw[DEplava!55] (-0.23,-4.44)--(-.10,-4.44);
\node[fill=white,inner sep=.2pt,text=DEplava] at (-0.23,-4.44) {60};
\draw (-.10,-4.44)--(2.55,-4.44);
\draw[fill=white] (2.55,-4.464)--(2.63,-4.464) arc(-90:90:.024)--(2.55,-4.416)--cycle;
\draw (2.654,-4.44)--(2.73,-4.44)--(2.73,-4.44)--(2.87,-4.44);
\draw[DEplava!55] (-0.61,-4.5)--(-.10,-4.5);
\node[fill=white,inner sep=.2pt,text=DEplava] at (-0.61,-4.5) {61};
\draw (-.10,-4.5)--(2.55,-4.5);
\draw[fill=white] (2.55,-4.524)--(2.63,-4.524) arc(-90:90:.024)--(2.55,-4.476)--cycle;
\draw (2.654,-4.5)--(2.73,-4.5)--(2.73,-4.48)--(2.87,-4.48);
\draw[DEplava!55] (-0.99,-4.5600000000000005)--(-.10,-4.5600000000000005);
\node[fill=white,inner sep=.2pt,text=DEplava] at (-0.99,-4.5600000000000005) {62};
\draw (-.10,-4.5600000000000005)--(2.55,-4.5600000000000005);
\draw[fill=white] (2.55,-4.5840000000000005)--(2.63,-4.5840000000000005) arc(-90:90:.024)--(2.55,-4.5360000000000005)--cycle;
\draw (2.654,-4.5600000000000005)--(2.73,-4.5600000000000005)--(2.73,-4.5200000000000005)--(2.87,-4.5200000000000005);
\draw[DEplava!55] (-1.37,-4.62)--(-.10,-4.62);
\node[fill=white,inner sep=.2pt,text=DEplava] at (-1.37,-4.62) {63};
\draw (-.10,-4.62)--(2.55,-4.62);
\draw[fill=white] (2.55,-4.644)--(2.63,-4.644) arc(-90:90:.024)--(2.55,-4.596)--cycle;
\draw (2.654,-4.62)--(2.73,-4.62)--(2.73,-4.5600000000000005)--(2.87,-4.5600000000000005);
\draw[fill=white] (2.82,-4.62) .. controls (3.02,-4.62) and (3.16,-4.590000000000001) .. (3.28,-4.44) .. controls (3.16,-4.29) and (3.02,-4.260000000000001) .. (2.82,-4.260000000000001) .. controls (2.94,-4.4) and (2.94,-4.48) .. cycle;
\draw (3.28,-4.44)--(3.38,-4.44);
\draw[fill=white] (3.38,-4.53)--(3.61,-4.44)--(3.38,-4.3500000000000005)--cycle;\draw[fill=white] (3.64,-4.44) circle(.03);
\draw (3.67,-4.44)--(4.10,-4.44);
\node[anchor=west,inner sep=1pt] at (4.1,-4.44) {O8};
\node[anchor=south,text=DEplava,inner sep=0pt] at (3.95,-4.420000000000001) {(12)};
\draw (-2.25,-4.71)--(-2.02,-4.71);
\draw[fill=white] (-2.02,-4.8)--(-1.80,-4.71)--(-2.02,-4.62)--cycle;
\draw[fill=white] (-1.80,-4.765) circle(.025);
\draw (-1.84,-4.655)--(2.24,-4.655);
\draw (-1.775,-4.765)--(2.32,-4.765);
\fill (2.24,-4.655) circle(.012);\fill (2.32,-4.765) circle(.012);
\node[anchor=east,inner sep=1pt] at (-2.25,-4.71) {I9};
\node[anchor=south,text=DEplava,inner sep=0pt] at (-2.30,-4.65) {(9)};
\draw[fill=white] (3.63,-4.8)--(3.40,-4.71)--(3.63,-4.62)--cycle;
\draw[fill=white] (3.40,-4.765) circle(.025);
\draw (3.44,-4.655)--(2.4,-4.655);
\draw (3.375,-4.765)--(2.48,-4.765);
\fill (2.4,-4.655) circle(.012);\fill (2.48,-4.765) circle(.012);
\draw (3.63,-4.71)--(4.1,-4.71);\node[anchor=west] at (4.1,-4.71) {I10};\node[anchor=south,text=DEplava,inner sep=0pt] at (3.95,-4.69) {(11)};
\draw[DEplava,decorate,decoration={brace,amplitude=2pt}] (-2.75,-4.82)--(-2.75,1.60) node[midway,left=3pt]{ulazi};
\draw[DEcrvena] (1.20,-4.41) ellipse (1.37cm and .30cm);
\node[anchor=north,text=DEtekst] (fuses) at (1.2,-5.05) {Programabilne veze};\draw[DEplava,->] (fuses.north)--(1.2,-4.73);
\coordinate (palout1) at (4.78,-0.24);
\coordinate (palout2) at (4.78,-0.84);
\coordinate (palout3) at (4.78,-1.44);
\coordinate (palout4) at (4.78,-2.04);
\coordinate (palout5) at (4.78,-2.6399999999999997);
\coordinate (palout6) at (4.78,-3.24);
\coordinate (palout7) at (4.78,-3.84);
\coordinate (palout8) at (4.78,-4.44);
\coordinate (palin10) at (4.78,-4.71);
\end{tikzpicture}

\column{.37\textwidth}\fontsize{9}{10}\selectfont
Maksimalno 10+\textcolor{DEcrvena}{6} ulaza odnosno 2+\textcolor{DEcrvena}{6} izlaza\par
„Bilo koja“ kombinaciona mreža koja ima maksimalno 10+X ulaza, odnosno 2+6-X izlaza\\X = 0,1,..6\par
20-pinsko pakovanje\\10. Pin GND\\20. Pin VDD\par
\tikz[remember picture,baseline]{\node[anchor=base west,inner sep=0pt] (palann1) {izlaz};}\par
\tikz[remember picture,baseline]{\node[anchor=base west,inner sep=0pt] (palannio) {Može biti i izlaz i ulaz};}\par
Ukupan broj programabilnih veza u ovoj komponenti je 2048 u I matrici. ILI matrica je fiksna. Po svakom I kolu je 32 moguća člana u proizvodu. Po svakom ILI kolu je moguće sumirati 7 proizvoda. I ovo je za današnju tehnologiju jednostavna programabilna komponenta.\par
\tikz[remember picture,baseline]{\node[anchor=base west,inner sep=0pt] (palann8) {izlaz};}\\\tikz[remember picture,baseline]{\node[anchor=base west,inner sep=0pt] (palann10) {ulaz};}
\begin{tikzpicture}[remember picture,overlay,DEplava,->,line width=.3pt]
\draw (palann1.west)--(palout1);
\draw (palann8.west)--(palout8);
\draw (palann10.west)--(palin10);
\draw (palannio.west)--(palout2);
\draw (palannio.west)--(palout3);
\draw (palannio.west)--(palout4);
\draw (palannio.west)--(palout5);
\draw (palannio.west)--(palout6);
\draw (palannio.west)--(palout7);
\end{tikzpicture}
\end{columns}

\end{frame}
